\documentclass[11pt]{article}

\usepackage{amsmath}
\usepackage{amssymb}
\usepackage{bm}
\usepackage{geometry}
\geometry{margin=1in}

\title{Equations and Variable Descriptions\\
\large Hearing the Room Through the Shape of the Drum}
\author{}
\date{}

\begin{document}

\maketitle

\section{Surface Vibration Model}

\subsection{General wave equation}

\begin{equation}
\rho \frac{\partial^2 u}{\partial t^2}
+
c(\bm{x}) \frac{\partial u}{\partial t}
-
T \nabla^2 u
+
D \nabla^4 u
=
f(\bm{x},t)
\end{equation}

\textbf{Description of terms:}

\begin{itemize}
    \item $u(\bm{x},t)$: out-of-plane displacement of the surface at spatial location $\bm{x}$ and time $t$.
    \item $\rho$: areal mass density of the surface.
    \item $c(\bm{x})$: spatially dependent damping coefficient.
    \item $T$: in-plane tension of the surface, dominant for membranes.
    \item $D$: flexural rigidity, dominant for thin plates.
    \item $\nabla^2 u$: Laplacian of the displacement, related to membrane-like wave behavior.
    \item $\nabla^4 u$: biharmonic operator, related to bending behavior of plates.
    \item $f(\bm{x},t)$: external forcing per unit area, induced by acoustic pressure.
\end{itemize}

For membranes:

\begin{equation}
D = 0
\end{equation}

For thin plates under negligible tension:

\begin{equation}
T = 0
\end{equation}

\section{Modal Decomposition}

\subsection{Expansion into vibration modes}

\begin{equation}
u(\bm{x},t)
=
\sum_{k=1}^{K}
\phi_k(\bm{x}) q_k(t)
\end{equation}

\textbf{Description of terms:}

\begin{itemize}
    \item $\phi_k(\bm{x})$: spatial eigenfunction, also called the $k$-th mode shape.
    \item $q_k(t)$: temporal modal coordinate of the $k$-th mode.
    \item $K$: number of modes used in the approximation.
\end{itemize}

\subsection{Second-order modal oscillator}

\begin{equation}
\ddot{q}_k(t)
+
2\zeta_k \omega_k \dot{q}_k(t)
+
\omega_k^2 q_k(t)
=
\alpha_k p(t)
+
\eta_k(t)
\end{equation}

\textbf{Description of terms:}

\begin{itemize}
    \item $q_k(t)$: temporal activation of the $k$-th mode.
    \item $\dot{q}_k(t)$: first temporal derivative, modal velocity.
    \item $\ddot{q}_k(t)$: second temporal derivative, modal acceleration.
    \item $\omega_k$: natural angular frequency of the $k$-th mode.
    \item $\zeta_k$: damping ratio of the $k$-th mode.
    \item $\alpha_k$: coupling coefficient between acoustic pressure and the $k$-th mode.
    \item $p(t)$: acoustic pressure signal driving the surface.
    \item $\eta_k(t)$: noise or modeling error for the $k$-th mode.
\end{itemize}

\section{Frequency-Domain Modal Transfer Function}

\subsection{Modal response in the frequency domain}

\begin{equation}
Q_k(\omega)
=
G_k(\omega) P(\omega)
\end{equation}

\textbf{Description of terms:}

\begin{itemize}
    \item $Q_k(\omega)$: Fourier transform of the modal coordinate $q_k(t)$.
    \item $P(\omega)$: Fourier transform of the acoustic pressure $p(t)$.
    \item $G_k(\omega)$: transfer function of the $k$-th vibration mode.
    \item $\omega$: angular frequency.
\end{itemize}

\subsection{Second-order oscillator transfer function}

\begin{equation}
G_k(\omega)
=
\frac{\alpha_k}
{-\omega^2
+
j 2 \zeta_k \omega_k \omega
+
\omega_k^2}
\end{equation}

\textbf{Description of terms:}

\begin{itemize}
    \item $G_k(\omega)$: frequency response of the $k$-th mode.
    \item $\alpha_k$: modal coupling coefficient.
    \item $j$: imaginary unit.
    \item $\zeta_k$: damping ratio.
    \item $\omega_k$: natural angular frequency of the mode.
    \item $\omega$: evaluated angular frequency.
\end{itemize}

\section{Surface Displacement as a Modal Convolution}

\begin{equation}
u(\bm{x},t)
=
\sum_{k=1}^{K}
\phi_k(\bm{x})
\left(
p(t) * g_k(t)
\right)
\end{equation}

\textbf{Description of terms:}

\begin{itemize}
    \item $g_k(t)$: impulse response of the $k$-th mode.
    \item $*$: convolution operator.
    \item $p(t) * g_k(t)$: acoustic pressure filtered by the mechanical response of the $k$-th mode.
\end{itemize}

\section{Sound Pressure and Source Signal}

\begin{equation}
p(t)
=
\gamma s(t)
\end{equation}

\textbf{Description of terms:}

\begin{itemize}
    \item $s(t) \in [-1,1]$: dimensionless source audio signal.
    \item $p(t)$: physical acoustic pressure in Pascals.
    \item $\gamma$: scale factor accounting for source amplitude and acoustic attenuation.
\end{itemize}

\section{Speckle-Based Vibration Measurements}

\subsection{Speckle shifts as surface gradients}

\begin{equation}
\bm{v}(\bm{x}_n,t)
=
\beta \nabla_{\bm{x}} u(\bm{x}_n,t)
\end{equation}

\textbf{Description of terms:}

\begin{itemize}
    \item $\bm{v}(\bm{x}_n,t) \in \mathbb{R}^2$: measured two-axis speckle shift at surface point $\bm{x}_n$.
    \item $\bm{x}_n$: the $n$-th measured surface point.
    \item $\beta$: optical transfer scale factor.
    \item $\nabla_{\bm{x}} u(\bm{x}_n,t)$: spatial gradient of the surface displacement.
\end{itemize}

\subsection{Forward model from sound to measured speckle shifts}

\begin{equation}
\bm{v}(\bm{x}_n,t)
\approx
\gamma \beta
\sum_{k=1}^{K}
\nabla \phi_k(\bm{x}_n)
\left(
s(t) * g_k(t)
\right)
+
\bm{\eta}(\bm{x}_n,t)
\end{equation}

\textbf{Description of terms:}

\begin{itemize}
    \item $\bm{v}(\bm{x}_n,t)$: measured two-axis vibration signal at point $\bm{x}_n$.
    \item $\nabla \phi_k(\bm{x}_n)$: spatial gradient of the $k$-th mode shape at point $\bm{x}_n$.
    \item $s(t)$: latent source sound signal to be recovered.
    \item $g_k(t)$: impulse response of the $k$-th mode.
    \item $\bm{\eta}(\bm{x}_n,t)$: measurement noise or modeling error.
    \item $\gamma \beta$: combined acoustic and optical scale factor.
\end{itemize}

\section{Fourier Transform of the Measured Vibrations}

\begin{equation}
\bm{V}(\bm{x}_n,\omega)
\equiv
\mathcal{F}
\left\{
\bm{v}(\bm{x}_n,t)
\right\},
\qquad
\bm{V}(\bm{x}_n,\omega)
\in
\mathbb{C}^2
\end{equation}

\textbf{Description of terms:}

\begin{itemize}
    \item $\mathcal{F}\{\cdot\}$: Fourier transform operator.
    \item $\bm{V}(\bm{x}_n,\omega)$: complex frequency-domain representation of the two-axis vibration signal.
    \item $\mathbb{C}^2$: two-dimensional complex vector space.
\end{itemize}

\section{Mode Frequency Estimation}

\begin{equation}
\left\{
\hat{\omega}_k
\right\}_{k=1}^{K}
\leftarrow
\operatorname{robustfindpeaks}
\left(
\left|
\bm{V}(\bm{x}_n,\omega)
\right|
\right)
\end{equation}

\textbf{Description of terms:}

\begin{itemize}
    \item $\hat{\omega}_k$: estimated natural frequency of the $k$-th mode.
    \item $\operatorname{robustfindpeaks}(\cdot)$: robust peak detection procedure used to identify modal frequencies.
    \item $|\bm{V}(\bm{x}_n,\omega)|$: magnitude spectrum of the measured vibration signal.
\end{itemize}

\section{Mode Shape Gradient Estimation}

\begin{equation}
\nabla \hat{\phi}_k(\bm{x}_n)
=
\operatorname{Re}
\left\{
\frac{
\bm{V}(\bm{x}_n,\hat{\omega}_k)
\cdot
V_1(\bm{x}_0,\hat{\omega}_k)^*
}{
\mathbb{E}_{n,a}
\left[
\left|
\bm{V}(\bm{x}_n,\hat{\omega}_k)
\right|
\right]
\cdot
\left|
V_1(\bm{x}_0,\hat{\omega}_k)
\right|
}
\right\}
\end{equation}

\textbf{Description of terms:}

\begin{itemize}
    \item $\nabla \hat{\phi}_k(\bm{x}_n)$: estimated gradient of the $k$-th mode shape at point $\bm{x}_n$.
    \item $\operatorname{Re}\{\cdot\}$: real part operator.
    \item $\bm{V}(\bm{x}_n,\hat{\omega}_k)$: measured vibration spectrum at the estimated modal frequency.
    \item $V_1(\bm{x}_0,\hat{\omega}_k)$: first-axis vibration measurement at reference point $\bm{x}_0$.
    \item $(\cdot)^*$: complex conjugate.
    \item $\mathbb{E}_{n,a}[\cdot]$: average over all points $n$ and vibration axes $a$.
\end{itemize}

\section{Sound Recovery Optimization}

\begin{equation}
\hat{s}^{\mathrm{inv}}(t)
=
\arg\min_{s(t),\alpha_k}
\left\|
\bm{v}(\bm{x}_n,t)
-
\sum_{k=1}^{K}
\nabla \hat{\phi}_k(\bm{x}_n)
\left(
s(t) * \hat{g}_k(t)
\right)
\right\|_2^2
+
\lambda
\left\|
\dot{s}(t)
\right\|_2^2
\end{equation}

\textbf{Description of terms:}

\begin{itemize}
    \item $\hat{s}^{\mathrm{inv}}(t)$: recovered source sound signal.
    \item $s(t)$: latent sound signal optimized during recovery.
    \item $\alpha_k$: modal coupling coefficient optimized jointly with the signal.
    \item $\nabla \hat{\phi}_k(\bm{x}_n)$: estimated gradient of the $k$-th mode shape.
    \item $\hat{g}_k(t)$: estimated impulse response of the $k$-th mode.
    \item $\lambda$: regularization weight.
    \item $\dot{s}(t)$: temporal derivative of the recovered sound signal.
    \item $\|\cdot\|_2^2$: squared Euclidean norm.
\end{itemize}

\subsection{Impulse response from the transfer function}

\begin{equation}
\hat{g}_k(t)
\leftarrow
\operatorname{iFFT}
\left(
G_k(\omega)
\right)
\end{equation}

\textbf{Description of terms:}

\begin{itemize}
    \item $\operatorname{iFFT}(\cdot)$: inverse Fast Fourier Transform.
    \item $G_k(\omega)$: analytical modal transfer function.
    \item $\hat{g}_k(t)$: estimated time-domain modal impulse response.
\end{itemize}

\section{Robust Mode Estimation}

\subsection{Spectral standard deviation across surface points}

\begin{equation}
\sigma(\omega)
=
\operatorname{std}_n
\left(
\left|
\bm{V}(\bm{x}_n,\omega)
\right|
\right)
\end{equation}

\textbf{Description of terms:}

\begin{itemize}
    \item $\sigma(\omega)$: standard deviation of spectral magnitudes across surface points.
    \item $\operatorname{std}_n(\cdot)$: standard deviation over measured points $n$.
    \item $|\bm{V}(\bm{x}_n,\omega)|$: magnitude spectrum of the vibration signal at point $\bm{x}_n$.
\end{itemize}

\subsection{Total variation of a mode-shape gradient}

\begin{equation}
\operatorname{TV}
\left(
\nabla \phi_k
\right)
=
\sum_{x,y}
\left\|
\frac{\partial}{\partial x}
\nabla \phi_k(x,y)
\right\|_2
+
\left\|
\frac{\partial}{\partial y}
\nabla \phi_k(x,y)
\right\|_2
\end{equation}

\textbf{Description of terms:}

\begin{itemize}
    \item $\operatorname{TV}(\nabla \phi_k)$: total variation of the gradient field of the $k$-th mode shape.
    \item $\frac{\partial}{\partial x}\nabla \phi_k(x,y)$: spatial derivative of the mode-gradient field along the $x$ direction.
    \item $\frac{\partial}{\partial y}\nabla \phi_k(x,y)$: spatial derivative of the mode-gradient field along the $y$ direction.
    \item $\|\cdot\|_2$: Euclidean norm.
\end{itemize}

\section{Calibration-Based Baseline}

\subsection{Reference convolution model}

\begin{equation}
v_a^{\mathrm{ref}}(\bm{x}_n,t)
\approx
s^{\mathrm{ref}}(t)
*
h_{n,a}(t)
\end{equation}

\textbf{Description of terms:}

\begin{itemize}
    \item $v_a^{\mathrm{ref}}(\bm{x}_n,t)$: reference vibration signal at point $\bm{x}_n$ and axis $a$.
    \item $s^{\mathrm{ref}}(t)$: known reference audio signal used for calibration.
    \item $h_{n,a}(t)$: impulse response from the reference sound to point $\bm{x}_n$ and vibration axis $a$.
    \item $a \in \{1,2\}$: vibration axis index.
\end{itemize}

\subsection{Inverse filter estimation}

\begin{equation}
\hat{h}_{n,a}^{-1}(t)
=
\arg\min_{h(t)}
\left\|
v_a^{\mathrm{ref}}(\bm{x}_n,t)
*
h(t)
-
s^{\mathrm{ref}}(t)
\right\|_2^2
\end{equation}

\textbf{Description of terms:}

\begin{itemize}
    \item $\hat{h}_{n,a}^{-1}(t)$: estimated inverse filter for point $\bm{x}_n$ and axis $a$.
    \item $h(t)$: candidate inverse filter optimized by least squares.
    \item $v_a^{\mathrm{ref}}(\bm{x}_n,t) * h(t)$: calibrated reconstruction of the reference signal.
\end{itemize}

\subsection{Recovered signal using calibrated inverse filters}

\begin{equation}
\hat{s}^{\mathrm{new}}(t)
=
\frac{1}{2N}
\sum_{a=1}^{2}
\sum_{n=1}^{N}
v_a^{\mathrm{new}}(\bm{x}_n,t)
*
\hat{h}_{n,a}^{-1}(t)
\end{equation}

\textbf{Description of terms:}

\begin{itemize}
    \item $\hat{s}^{\mathrm{new}}(t)$: recovered signal for a new measurement.
    \item $v_a^{\mathrm{new}}(\bm{x}_n,t)$: new measured vibration signal at point $\bm{x}_n$ and axis $a$.
    \item $\hat{h}_{n,a}^{-1}(t)$: calibrated inverse filter.
    \item $N$: number of measured surface points.
    \item $2N$: total number of measurement channels, since each point has two vibration axes.
\end{itemize}

\section{Implementation Constants}

\begin{equation}
\zeta_k = 0.01
\end{equation}

\textbf{Description of term:}

\begin{itemize}
    \item $\zeta_k$: damping ratio fixed to $0.01$ in the experiments.
\end{itemize}

\begin{equation}
\lambda = 1
\end{equation}

\textbf{Description of term:}

\begin{itemize}
    \item $\lambda$: regularization weight used in the sound recovery optimization.
\end{itemize}

\begin{equation}
T_{\mathrm{conv}} = 4096
\end{equation}

\textbf{Description of term:}

\begin{itemize}
    \item $T_{\mathrm{conv}}$: number of taps used for the inverse filters in the calibration baseline.
\end{itemize}

\end{document}